\documentclass[11pt]{article}
\usepackage[T1]{fontenc}
\usepackage{lmodern,amsmath,amssymb,amsthm,mathtools}
\usepackage[margin=0.95in]{geometry}
\usepackage{microtype,booktabs,longtable,array,enumitem,xcolor,graphicx}
\usepackage{fancyhdr}
\usepackage[colorlinks=true,linkcolor=blue!45!black,urlcolor=blue!55!black,citecolor=blue!45!black]{hyperref}
\hypersetup{pdftitle={Smaller explicit Diophantine forms of ZFC consistency},
pdfauthor={AI-assisted research prepared for Vladimir Reshetnikov},
pdfsubject={Explicit equations, proofs, coefficient recipes, and reproducible audits},
pdfkeywords={Diophantine equations, ZFC consistency, K215, Jones reduction}}
\usepackage{listings}
\lstset{basicstyle=\ttfamily\small,breaklines=true,columns=fullflexible,keepspaces=true}
\setlength{\parskip}{0.35em}
\setlength{\parindent}{0pt}
\setlength{\emergencystretch}{2em}
\setlist{itemsep=0.2em,topsep=0.3em}
\pagestyle{fancy}\fancyhf{}
\fancyhead[L]{\small Explicit Diophantine forms of ZFC consistency}
\fancyhead[R]{\small Research synthesis}
\fancyfoot[C]{\thepage}
\setlength{\headheight}{14pt}
\newtheorem{theorem}{Theorem}[section]
\newtheorem{lemma}[theorem]{Lemma}
\newtheorem{proposition}[theorem]{Proposition}
\newtheorem{corollary}[theorem]{Corollary}
\theoremstyle{definition}\newtheorem{definition}[theorem]{Definition}
\theoremstyle{remark}\newtheorem{remark}[theorem]{Remark}
\newcommand{\N}{\mathbb N}\newcommand{\Z}{\mathbb Z}\newcommand{\Q}{\mathbb Q}
\newcommand{\NN}{\mathbb N}\newcommand{\ZZ}{\mathbb Z}\newcommand{\QQ}{\mathbb Q}
\newcommand{\ZFC}{\mathrm{ZFC}}\newcommand{\Con}{\operatorname{Con}}
\newcommand{\Sq}{\operatorname{Sq}}\newcommand{\pc}{\operatorname{pc}}
\newcommand{\ord}{\operatorname{ord}}\newcommand{\lcm}{\operatorname{lcm}}
\newcommand{\pin}{\texttt{8c6eec9d19ef80538c8e61ff6e0180040583446a}}
\newcommand{\repo}{\url{https://github.com/VladimirReshetnikov/ProveIt}}
\begin{document}
\begin{center}
{\LARGE\bfseries Smaller explicit Diophantine forms\par of ZFC consistency}\\[0.7em]
{\large Mixed radices, cofinal bases, and norm cancellations}\\[1em]
Prepared for Vladimir Reshetnikov\\[0.2em]
AI-assisted research report, 22 October 2026
\end{center}
\begin{abstract}
We compare the seven current ProveIt constructions, numbered 24--30,
and combine features that their separate presentations leave uncombined.
The main new variable reduction gives one explicitly specified polynomial
in sixteen unrestricted integers. Its degree is
$813056\cdot5^{32}+10024$, and its delivered arithmetic circuit uses
1865 binary operations. A seventeen-integer alternative has degree
$406528\cdot5^{32}+27512$ and 1874 operations. A cofinal family of square
radices removes one positivity witness; a smaller signed separating weight
also lowers the degree. For positive unknowns we obtain a twelve-variable
equation of degree 1065, and separately reduce the 1081-degree equation
from 111 to 108 operations. Three coordinate shifts give a 31-variable
quadratic system using 62 operations, or a quartic equation using 127.
All families can use a new, fully materialized four-coefficient tuple with
$B=2^{1064921}$, a substantial reduction from $2^{1953125}$.
We give ordinary mathematical proofs, literal arithmetic circuits,
coefficient generators, and independent finite audits. The equivalences
are formalizable in ZFC; the new reductions are not claimed to be new
end-to-end Lean or Rocq certificates, globally minimal presentations,
or proofs of ZFC consistency.
\end{abstract}
\newpage
\tableofcontents
\newpage

\section{Outcome, comparison, and interpretation}
\label{sec:outcome}

Write $\N=\{0,1,2,\ldots\}$ and $\N_{>0}=\{1,2,\ldots\}$.
For every construction delivered here, with its explicitly stated domain
$D$, the intended theorem is
\begin{equation}\label{eq:main-theorem}
 \ZFC\vdash
 \bigl(\neg\exists\boldsymbol{x}\in D:
           P(\boldsymbol{x})=0\bigr)
 \longleftrightarrow\Con(\ZFC).
\end{equation}
For a system, $P=0$ in this display means that all listed equations hold.
The consistency predicate means that no finite derivation of contradiction
exists in the fixed effective presentation of ZFC used by the K215 endpoint.
The proof establishes the equivalent positive statement: a solution exists
if and only if a finite ZFC refutation exists.

There is no single ordering that makes an equation with fewer variables,
another with smaller degree, and a third with fewer arithmetic operations
all ``simplest.'' We therefore retain the useful new points on these
tradeoffs and identify their exact gains. An operation is one binary
$+$, $-$, or multiplication, including multiplication by a fixed numeral.
Reuse and equality comparisons are free. Construction of fixed integers
is accounted for separately from evaluation of a polynomial at an input.
The degree counts quantified variables only: coefficient ports have degree
zero. These conventions agree with the frozen repository catalogue.

\begin{center}\small
\begin{tabular}{@{}p{0.25\textwidth}p{0.32\textwidth}p{0.32\textwidth}@{}}
\toprule
Criterion & Frozen predecessor & This report\\
\midrule
Unrestricted integer count & 17 variables, 20607 operations &
16 variables, 1865 operations\\
17-integer alternative & Degree $424296\cdot5^{33}+53824$ &
Degree $406528\cdot5^{32}+27512$\\
Single-equation degree in 12 positive variables & Degree 1069, 123 operations &
Degree 1065, 118 operations\\
Single-equation cost in 12 positive variables & Degree 1081, 111 operations &
Degree 1081, 108 operations\\
Quadratic system in 31 positive variables & 22 equations, 65 operations &
22 equations, 62 operations\\
31-positive-variable quartic & 130 operations & 127 operations\\
Smallest materialized base & $B=2^{1953125}$ & $B=2^{1064921}$\\
Four coefficient payloads & 977220 raw bytes & 533099 raw bytes\\
\bottomrule
\end{tabular}
\end{center}

The sixteen-integer polynomial has degree
$813056\cdot5^{32}+10024$, which is also lower than the predecessor's
seventeen-integer degree. It is still enormous. For a readable low-degree
presentation, the 22-equation quadratic system is the more useful choice.
For a compact direct positive-domain circuit, the 108-operation equation
is preferable. For the fewest positive variables, the existing Jones
nine-variable route remains relevant: this report does not improve the
count nine. The established 67-operation single equation and
59-operation higher-degree systems are also not beaten in their own
operation-count categories.

The comparisons concern explicit constructions in the inspected repository.
They are not assertions of universal optimality. In particular, Sun's
September 2026 preprint states an effective ten-integer representation
theorem for every recursively enumerable set \cite{sun10}. That theorem
gives a smaller general existence bound than sixteen. Its different
source coding has not been instantiated into a fixed K215 polynomial in
this delivery; calling the sixteen-integer circuit the globally smallest
known integer equation would therefore be unjustified.

\subsection{Sources and the frozen snapshot}

The source repository is \repo. All comparisons and source copies in this
archive use commit\par
\begin{center}\small\pin\end{center}
in the directory
\begin{quote}\small\ttfamily
SetTheory/Cardinals/docs/reports/hilbert-tenth-problem/\\
con-zfc-diophantine-forms/
\end{quote}
The main catalogue at that commit
already incorporates its audit of the fourth round of submissions.
We compare those audited files, rather than treating a historical
summary or a package's claim about its predecessors as current evidence.

The useful contributions are complementary. Package 24 provides a tight
separating weight and a padding argument for general odd radix dimensions.
Package 25 provides mixed-radix local decoding, compact clause bands, and
population adjustment in fields that are identically zero. Package 30
provides a trace/product norm evaluator and the 31-variable quartic.
Package 29 provides a signed Jones finalizer with a stack of positivity
certificates. Combining the last two is especially effective: the signed
construction had not yet used the smaller quartic or the compact evaluator.

\section{The proof interfaces and what is verified}
\label{sec:interfaces}

The inherited proofs used here have four precise roles.
\begin{enumerate}
\item The literal 215-state machine on the blank tape halts if and only if
ZFC has a finite refutation. The pinned Lean endpoint states this as
\texttt{KEndpoint215.blank\_halts\_iff\_zfc\_refutation}; its statement is
included in the archive. We use the corresponding ordinary mathematical
calibration, rather than an assumption that ZFC proves the soundness of
Lean as a whole.
\item The fixed 943-tile language represents finite blank computations.
A periodic tiling with a Start marker and no Forbidden tile yields a
halting computation. Conversely a finite halt can be padded to arbitrarily
large independent widths and heights. These are the exact semantic
properties needed when we change the numerical layout.
\item The native-window and Pell compiler translates that language into
the direct polynomial interface. Its soundness is for arbitrary
solutions, and its completeness constructs positive witnesses. The
canonical witnesses additionally have square scale $w$, odd scale $s$,
and an odd strong-Pell multiplier where the signed reductions use them.
The materialized tuple must satisfy this whole interface; modular
identities of four unrelated integers would not suffice.
\item The Jones reduction translates a normalized quartic into its
eight-coordinate ratio predicate, or into a nine-variable equation over
$\N$. The signed finalizer uses that particular predicate and its
cofinal-radix completeness proof. It is not a generic conversion of an
arbitrary nine-variable polynomial into sixteen integer variables.
\end{enumerate}

The new coefficient adapter is proved in Section~\ref{sec:coefficients}.
The new arithmetic reductions and the signed construction are proved in
the later sections. These proofs, together with the cited inherited
lemmas, give \eqref{eq:main-theorem} by ordinary mathematics formalizable
in ZFC. The archive includes the relevant inherited proof notes, literal
source data, and source locations so that this dependency is reviewable.
The native compiler and new finalizers remain outside the existing
end-to-end proof-assistant certificates.

Finite checking has a distinct, narrower role. It verifies the literal
tile list, numerical coefficients, convolution identities, size bounds,
operation counts, coordinate identities, norm identities, and degree
certificates. It does not search for a counterexample to consistency,
simulate an unbounded machine run, or establish the universal compiler
theorem by testing finitely many cases.

\section{A smaller complete coefficient tuple}
\label{sec:coefficients}

The new backend combines the mixed-radix local decoder and compact bands
of package 25 with the general odd-factor Boolean padding argument of
package 24. Neither original numerical tuple is reused piecemeal.
All four ports are constructed together and then audited together.

\subsection{The finite language and its literal encoding}

The file \texttt{coefficients/k215.tm} contains the 430 literal actions
of the 215-state machine. Its SHA-256 is
\begin{center}\scriptsize\ttfamily
fdbcef0825886035c3091c8537a69535bc96f9e265f0ae45863543d40ae70090.
\end{center}
The generator never runs this machine. It reads the finite transition
table, constructs the same 943 ordered tiles as package 25, and checks
the canonical tile-list hash. A tile has west, north, east, and south
edge labels. The language contains 430 transition tiles, two halted-head
loops, two unchanged-bit tiles, 398 right-arrival tiles, 90 left-arrival
tiles, four finish tiles, four initialization tiles, three reset tiles,
nine seam tiles, and one Forbidden tile. Start is placed first,
Forbidden second, and the remaining tiles are lexicographically ordered.

For clarity, the finite semantic contract retained from the source is as
follows. A marked periodic plane forces initialization between reset
boundaries. The directed signal colors update exactly one head at each
time step and cannot cross a boundary. A new reset row can occur only
after a halted head has appeared. Thus a component containing Start
contains a finite blank halting computation. Conversely, a finite halt
admits arbitrarily wide blank margins and arbitrarily many halted padding
rows; the seam tiles close these to a helical periodic tableau. Its width
and height can be chosen independently among all sufficiently large
integers, in particular among powers of the prime used below.

Let
\[
 b_c=121,\qquad L_c=8801,\qquad d_c=b_cL_c=1064921,
 \qquad \rho=2^{121},\qquad B=\rho^{8801}.
\]
The native support is
\begin{align*}
 \mathcal E={}&\{0,\ldots,942\}\cup\{983,\ldots,992\}
              \cup\{1003,\ldots,1032\}\\
 &{}\cup\{1063,1064,1095,1096,2195,3296\}.
\end{align*}
It has 989 elements. The two dummy positions are 1063 and 1064;
the first permits the harmless binary variation used for index control.
For the four edge slots use copy offsets
$o_0=1023,o_1=1013,o_2=1003,o_3=983$.
Put $w_j=2^{10+2j}$, for $0\le j<44$.
If tile $i$ has labels $x_{i,t}$, its clause coefficient is
\[
 c_i=1+\sum_{t=0}^3\sum_{j=0}^9
       \operatorname{bit}_j(x_{i,t}+512)w_{10t+j}
       +\sum_{j=40}^{43}w_j.
\]
At the copy position $o_t+j$, set $c_{o_t+j}=w_{10t+j}$.
At the anchors $1095,1096,2195,3296$, set the coefficients to
$w_{40},w_{41},w_{42},w_{43}$, respectively. Both dummy coefficients
are zero. The clause mask is
\[
 \mu=1022+\sum_{j=0}^{43}w_j.
\]
Its bit length is 97 and its population is 53.

The two clause targets are $T_1=4393,T_2=5495$ and the spatial
displacement is 42. Define the fixed integers
\begin{align*}
 D_C={}&\rho^{30}+\rho^{62}+\rho^{1100}+\rho^{2200}
       +\sum_{e\in\mathcal E}c_e
                      (\rho^{4393-e}+\rho^{5495-e}),\\
 D_R={}&\rho^{42},\\
 M_C={}&B-1-\sum_{\substack{e\in\mathcal E\\e\ne1}}\rho^e
                -2\rho^{1063},\\
 M_{F,\mathrm{native}}={}&4+\mu(\rho^{4393}+\rho^{5495})
      +\sum_{j=0}^9\rho^{1013+j}
      +\sum_{j=0}^9\rho^{1065+j}+\rho^{2195}+\rho^{3296}\\
 &{}+(2^{119}-1)\sum_{e=8792}^{8798}\rho^e
             +(2^{27}-1)\rho^{8799}.
\end{align*}
Every exponent is nonnegative. The four actual ports are
\begin{equation}\label{eq:new-ports}
 \boxed{\beta=B-1,\quad K_0=D_C+BD_R,\quad M_C,\quad
 M_F=M_{F,\mathrm{native}}+B-1.}
\end{equation}
In particular $M_F$ is already shifted. The unshifted mask is not the
fourth arithmetic input.

\subsection{Decoder, separation, and carries}

The low ten bits of a clause sum count the selected tiles. Before any
decoding there are at most 943 selectors, and all higher coefficients
are multiples of 1024. The mask 1022 therefore forces the selector count
to be either zero or one. Conditioned on this conclusion, each of the
44 remaining radix-four digits is a sum of two Boolean values. It has
no carry and its masked parity forces the two values to be equal.
This gives both the edge-copy constraints and the anchor constraints.
The order of these arguments matters: carry-freeness of the radix-four
digits is used only after the selector count has been recovered.

The overlap tests are unchanged. At position $1013+j$, the bare shift
30 pairs the south copy $983+j$ with the temporal north copy $1013+j$.
At position $1065+j$, the bare shift 62 pairs the east copy $1003+j$
with the spatial west copy $1023+j$ shifted by 42. The two anchor tests
at 2195 and 3296 pair copies at 1095 and 1096, respectively, with their
temporal copies. At a clause target, the reflected coefficient from
each support position is exactly $c_e$, and all other possible
contributions vanish.

These are identities on every native basis vector, not observations
about selected example tableaux. There are 69230 such finite identities
in the delivered audit. The two further support identities are
\begin{align*}
 \{(4393-e)\bmod8801:e\in\mathcal E\}
 &\cap\{(5495-e)\bmod8801:e\in\mathcal E\}=\varnothing,\\
 \{(2195-e)\bmod8801:e\in\mathcal E\}
 &\cap\{(3296-e)\bmod8801:e\in\mathcal E\}=\{0\}.
\end{align*}
The first preserves a clean clause band under arbitrary rotations;
the second forces whole-cell synchronization. The source's two-anchor
parity argument first excludes subdigit offsets, including the one
possible upper-dummy spill. The displayed intersection then excludes
all remaining nonzero offsets. Thus this part of soundness uses no
power-of-five assumption on the new dimensions.

Put
\[
 \mathcal M=(989+2)\left(2\sum_{e\in\mathcal E}c_e+6\right).
\]
Its bit length is 118, while that of $4\mathcal M+8$ is 120.
The complete conservative threshold is
\[
 \max(4\mathcal M+8,2\mu+4,2(2^{119}-1)+4,16)
       =2^{120}+2<\rho.
\]
This bound applies before selector extraction, including all allowed
dummy values. It preserves the no-carry and positive raw-slack
hypotheses of the source compiler.

The largest unshifted native product position is
$5495+3296=8791$. The padding masks occupy only positions 8792 through
8799, and $8799<8801$. They therefore test fields that are structurally
zero in every canonical tableau, even after dummy-bit changes. Their
population is $7\cdot119+27=860$. Consequently
\[
 \pc(M_C)=d_c-989,
 \qquad\pc(M_{F,\mathrm{native}})=2\cdot53+22+1+860=989.
\]
Both masks are strictly between zero and $B-1$, and
\[
 M_C\equiv2\pmod4,\qquad M_{F,\mathrm{native}}\equiv4\pmod8.
\]
The population identity and residues are verified on the materialized
integers as well as on their finite recipes.

\subsection{Completeness for the new odd dimensions}

It would be incorrect merely to use the new dimensions in the old
power-of-five completeness argument. Instead factor
\[
 d_c=11^2\cdot13\cdot677
\]
and choose the independently verified prime
$\lambda=2341$ and the step $k_c=2340=\lambda-1$.
The exact order data are
\begin{center}
\begin{tabular}{@{}lrrrr@{}}\toprule
Modulus & 121 & 13 & 677 & 2341\\\midrule
Order of $B$ & 10 & 12 & 52 & 60\\\bottomrule
\end{tabular}
\end{center}
Thus $\ord_{d_c}(B)=780$ and $k_c$ is a common multiple of all
orders. It is deliberately not claimed to be their least common multiple.
Also
\[
 B^{2340}\bmod2341^2=2556373=1+1092\cdot2341.
\]
The valuation $v_{2341}(B^{2340}-1)$ is therefore exactly one.
Writing $H_0=1+BD_C+B^2(D_R+1)$, the unit check is
\begin{center}
\begin{tabular}{@{}rrrrr@{}}\toprule
$p$ & $B\bmod p$ & $D_C\bmod p$ & $(B^2-1)\bmod p$ & $H_0\bmod p$\\\midrule
11 & 2 & 5 & 3 & 9\\
13 & 6 & 1 & 9 & 7\\
677 & 611 & 38 & 293 & 654\\
2341 & 2100 & 2051 & 1896 & 1225\\\bottomrule
\end{tabular}
\end{center}

Choose sufficiently large independent tableau width $w_c$ and height
$h_c$ that are powers of $\lambda$, and put $N_c=w_ch_c=\lambda^n$.
Require $N_c>d_c\lambda^2$ and $h_c>\lambda$; padding a fixed finite
halt permits these requirements. Put $q=B^{N_c}$. Since
$\lambda\equiv1$ modulo the relevant orders,
\[
 q\equiv B^{w_c}\equiv B\pmod{d_c\lambda}.
\]
Changing a free low dummy in cell $j$ changes the actual packed Pell
index by $-\Gamma B^j$, where
\[
 \Gamma=\rho^{1063}(q^2-1)
             [1+q(D_C+BD_R+B^{w_c})].
\]
Every prime factor of $d_cN_c$ is in the unit table. Thus
$\gcd(\Gamma,d_cN_c)=1$.

\begin{lemma}[Boolean coverage of the required modulus]
Put $T=N_c/\lambda$. Every residue modulo $d_cN_c$ is a sum of a
Boolean subset of
\[
 \{B^{2+k_cj}:0\le j<T\}\cup\{B^3\}.
\]
The chosen cells and their required spatial and temporal translates
can all be kept nonwrapping.
\end{lemma}
\begin{proof}
Let $C_0=d_c\lambda$. Then $d_cN_c=C_0T$ and $T>C_0$.
The valuation calculation and the binomial lifting identity show that
$B^{k_c}$ has order $T$ modulo $\lambda^n$ and order one modulo $d_c$.
Its first $T$ powers are exactly the $T$ residues $1+C_0z$
modulo $C_0T$. This follows either by the Chinese remainder theorem
and cardinality or directly from the order calculation.

Multiply the target by $B^{-2}$. Choose $\varepsilon\in\{0,1\}$
so that its difference from $\varepsilon B$ is nonzero modulo
$\lambda$. Its least residue $a$ modulo $C_0$ then satisfies
$1\le a<C_0<T$ and $\gcd(a,T)=1$.
Write the remaining target as $a+C_0u$. Solve
\[
 az_0+\frac{a(a-1)}2\equiv u\pmod T.
\]
The $a$ distinct orbit entries with additive coordinates
$z_0,z_0+1,\ldots,z_0+a-1$ have the desired sum. Include $B^3$
exactly when $\varepsilon=1$. All selections are Boolean.

The largest available cell position is
$2+(\lambda-1)(N_c/\lambda-1)=N_c-N_c/\lambda-\lambda+3$.
As $h_c>\lambda$, the width $w_c$ is less than $N_c/\lambda$.
Translations by one and by $w_c$ therefore remain below $N_c$.
The additional cell 3 is distinct and also safe.
\end{proof}

Start with a canonical tableau whose dummy bits are zero, and let its
packed index be $R_0$. Apply the lemma to
$(R_0-d_cw_c)\Gamma^{-1}$ modulo $d_cN_c$.
The modified tableau has $R\equiv d_cw_c\pmod{d_cN_c}$.
This is exactly the temporal rotation demanded by $X=2^R$.
The changes affect no selected tile, matched edge, Start marker, zero
padding field, or carry estimate. The inherited Pell completeness
argument now supplies all positive witnesses.

\subsection{Canonical parity and the signed interfaces}

Let $u_c=d_cN_c$, so $q=2^{u_c}$. All dimensions are odd. The
canonical Start digit and forbidden selector imply that the native
code $C$ is $1$ modulo 4. With
\[
 S=(C-1)+qF,\qquad
 T_m=(M_CJ+1)+q(M_{F,\mathrm{native}}J-1),\qquad M=q^2,
\]
the source packing identity is
\[
 R=(M-S)(M-1)+T_m=(M-1-S)M+(S+T_m).
\]
The native masks give disjoint binary support of $S,T_m$;
the exact population balance gives $\pc(T_m)=u_c+2$.
The upper complement then proves $\pc(R)=3u_c+2$.
The low residues also give $R\equiv3\pmod4$.

The canonical scale $w=2^{R-u_c}$ is consequently an integer square.
For the even main Pell parameter, its coordinate $c$ at the odd
rank $R$ is odd; choosing the strong rank $cR$ makes the normalized
multiplier odd, and its ordinary discriminant-scaled version is odd
as well. Finally, the inherited half-binomial valuation identity
$v_2(Y)=\pc(R)-2=3u_c$ shows that $s=Y/q^3$ is odd.
Thus the new tuple supplies the full canonical odd and square
properties, in addition to the unsigned compiler equivalence.

\begin{theorem}[Complete replacement of the four ports]
The tuple \eqref{eq:new-ports} satisfies the same blank-K215
native-window/Pell interface as the Wang943 source, including its
canonical odd and square witnesses. It can therefore be used as a
whole tuple in the compatible direct forms, the ordinary quartic,
and the Jones and signed constructions based on that quartic.
\end{theorem}

\subsection{Materialized size and reproducibility}

The archive contains the actual four integers, both as raw minimal
unsigned little-endian byte strings and as gzip files. Their sizes are
\begin{center}
\begin{tabular}{@{}lrrr@{}}\toprule
Port & Bit length & Raw bytes & Gzip bytes\\\midrule
$\beta$ & 1064921 & 133116 & 164\\
$K_0$ & 1070004 & 133751 & 9794\\
$M_C$ & 1064921 & 133116 & 264\\
$M_F$ & 1064922 & 133116 & 263\\\midrule
Total & & 533099 & 10485\\\bottomrule
\end{tabular}
\end{center}
The exponent of $B$ is about 45.48\% smaller than the preceding minimum;
the raw payload is about 45.45\% smaller. Compression is reported only
as a transport size, not as coefficient height. These are four actual
integers, fixed by a terminating finite recipe independently of whether
the machine halts. The much larger Jones index coefficients are separately
specified by exact finite recipes in the signed construction.

The builder and the independent audit use only Python's standard
library. The latter reconstructs the tuple without importing the builder
and checks every port against the delivered bytes. Both the generator
and audit can be rerun without network access.

\section{A 62-operation quadratic system and its quartic}
\label{sec:quartic}

For a presentation with small degree, the most readable new result is
a system of 22 quadratic equations in 31 strictly positive integers.
Three coordinated changes of coordinates reduce package 30's system
from 65 operations to 62. Its sum of squares is a single quartic using
127 operations, compared with 130 before the changes.

The 31 variables are
\begin{align*}
 &J,F,\alpha,z,f,h,i,o,j,s,w,\tau,\eta,\zeta,y,\gamma,\\
 &N,E,T,c,a,X,D,C_2,S,F_2,V,g,I,B_a,\Delta.
\end{align*}
All are strictly positive. Use the four fixed coefficients
\eqref{eq:new-ports} and the affine abbreviations
\[
 q=\beta J+1,\qquad k=\eta+\zeta,\qquad H=4a+3,
 \qquad C=q-F-\alpha.
\]
The equations are
\begin{align}
 S&=q^2, &N&=qS, &X&=wq,\label{eq:quad-start}\\
 a&=E+sN, &E&=X(a-E), &c&=k(a-E)+\eta,\\
 T&=E(c-\eta), &\tau^2&=T(T+k)+1,\\
 D&=X+ac+\gamma H, &\Delta&=a^2+H,\\
 C_2&=c^2, &D^2&=\Delta C_2+1,\\
 k&=jc-V+1+hE.&&
\end{align}
The two packing and transport equations are
\begin{align}
 jc-V&=(q-F)(N-q)-C(S-1)+(M_C+qM_F)J,\\
 (K_0+w)C&=F+z(q-1).
\end{align}
The remaining seven equations are
\begin{align}
 F_2&=f^2, &V&=of-c, &g&=y^2-V^2,\\
 I&=iC_2, &B_a&=I^2, &B_a&=\Delta(F_2-1),\\
 y^2-B_ag&=1.&&&\label{eq:quad-end}
\end{align}
Since every abbreviation is affine and every displayed product has
at most two affine factors, every residual has total degree at most two.
Writing the 22 residuals as $r_j$, the single polynomial
\begin{equation}\label{eq:new-quartic}
 Q_{31}=\sum_{j=1}^{22}r_j^2
\end{equation}
has degree exactly four: the coefficient of $\tau^4$ is one. Over the
stated positive domain, $Q_{31}=0$ is equivalent to all 22 equations.

\subsection{The three shifts and the positive zero sets}

The predecessor coordinates being changed are denoted by
$L_{\rm old},P_{\rm old},z_{\rm old}$. The exact substitutions are
\[
 L_{\rm old}=S-1,\qquad P_{\rm old}=T+k,
 \qquad z_{\rm old}=z+1.
\]
All other variables are untouched.

The old equations for $L_{\rm old}$ and $N$ become $S=q^2$ and
$N=qS$. In the packing equation $L_{\rm old}$ becomes $S-1$.
This removes two gates and introduces one, for a net saving of one.
The old first-root definition was
$P_{\rm old}=E(c-\eta)+k$; after the shift it is
$T=E(c-\eta)$. The first norm becomes $\tau^2=T(T+k)+1$.
This saves a second gate. Finally the transport identity
\[
 (K_0+w)C+(q-F)-(z+1)(q-1)-1
       =(K_0+w)C-F-z(q-1)
\]
saves a third gate when its two sides are compared directly.
These are polynomial identities. The sparse-ring audit verifies each of
the 22 residuals against both the predecessor's expression and its
literal sparse coefficient data after substitution.

Given a new positive solution, $S=q^2>1$ implies
$L_{\rm old}>0$, while $P_{\rm old}=T+k>0$ and
$z_{\rm old}=z+1>0$. The identities give a positive predecessor solution.
Conversely, on a positive predecessor solution,
$S=L_{\rm old}+1=q^2$, and
\[
 a-E=sN>0,\qquad c-\eta=k(a-E)>0,
 \qquad T=P_{\rm old}-k=E(c-\eta)>0.
\]
For the final shift, the inherited Pell/compiler theorem supplies
$C>0$, $X=wq=2^R$, $R>3q$, and $F<q$.
Consequently $w>q>F$, and the transport equation gives
\[
 (z_{\rm old}-1)(q-1)=(K_0+w)C-F>0.
\]
Thus $z_{\rm old}\ge2$, so the new $z$ is positive.
This positivity step is a use of the inherited arithmetic theorem,
not an inference from a finite test of a few tuples.
The transformations are inverse on the positive zero sets.

\begin{theorem}[The low-degree construction]
ZFC proves that the 22-equation system above has no positive integer
solution if and only if $\Con(\ZFC)$. Equivalently, ZFC proves the
same statement for the quartic $Q_{31}=0$.
\end{theorem}
\begin{proof}
The shifts give the proved bijection to package 30's positive quartic
system. The complete coefficient replacement theorem permits the new
tuple in that system. Its inherited Pell/compiler equivalence and
the literal K215 calibration give the theorem.
\end{proof}

\subsection{Cost and the quartic used by the Jones branch}

The system has 62 live binary gates: 33 multiplications and 29 additions
or subtractions. The sum-of-squares finalizer adds 22 residual
subtractions, 22 squares, and 21 additions, giving 127 gates.
Its expansion has 297 monomials when the four coefficient ports remain
independent formal symbols; the predecessor has 306 under that convention.
This count does not identify algebraic coefficient ports with quantified
unknowns.

Replacing every positive variable by its nonnegative successor gives a
normalized quartic: at all positive variables equal to one the residual
$E-X(a-E)$ equals one independently of the coefficient tuple. Thus the
shifted quartic is nonzero at the all-zero tuple, as required by Jones.
To keep the index coefficients and their exact source map simple, the
signed construction below uses the predecessor's literal 31-variable
quartic, rebound to the same new tuple. It has the same existential
meaning. Its separate source data are included, and are not confused
with the new 297-term quartic.

\section{Two improved twelve-positive-variable constructions}
\label{sec:positive-improved}

There are two different improvements.  The economical construction evaluates
exactly package 30's degree-$1081$ polynomial in its four coefficient ports
in $108$, instead of $111$,
binary arithmetic operations.  Its three-equation system costs $104$
operations, or $103$ after a justified shift of the positive transport
coordinate.  The second construction changes the polynomial: it has degree
$1065$ in the same twelve positive unknowns, compared with the previous
degree record $1069$.  It costs $118$ operations; its three-equation versions
cost $114$ and $113$ operations, with maximum residual degree $1056$.
The latter combines package 24's sharp weight, package 25's exclusion of
the negative main-root branch, and package 30's quadratic-remainder norm.
The smaller carrier permitted by the branch exclusion creates two additional
leading cancellations.  These are comparisons with the frozen repository
collection, not assertions of a global minimum.

All results in this section use one complete coefficient tuple satisfying
the native blank-compiler interface.  The fixed ports are
$\beta=B-1,K_0,M_C,M_F$, with $M_F$ the shifted high mask.  They are fixed
integers, not additional unknowns.  The coefficient construction elsewhere
in this report supplies the tuple.  Changing only one port is not justified.

\subsection{Common polynomial definitions}

The supplied variables are exactly
\[
 J,F,\alpha,z,h,i,T,s,w,\eta,\zeta,t\in\mathbb Z_{>0}.
\]
Every other displayed symbol denotes a polynomial expression in these
variables.  Put
\begin{align*}
 q&=\beta J+1,& C_w&=q-F-\alpha,&X&=wq,&Y&=sq^3,&E&=XY,\\
 k&=\eta+\zeta,&c&=kY+\eta,&a&=E+Y,&H&=4a+3,&\Delta&=a^2+H,\\
 \ell&=EkY,&A&=1+\ell(\ell+k),&G&=1+\Delta c^2,&b_0&=X+ac.
\end{align*}
In particular $\Delta=(a+2)^2-1$.  The unchanged index and transport
expressions are
\begin{align*}
 R&=((q-1)(q-F)+\alpha)(q^2-1)+(M_C+qM_F)J,\\
 N_k&=k-hE-R,\\
 N_t&=(K_0+w)C_w+q-F-z(q-1).
\end{align*}
Define the strong and auxiliary expressions
\begin{align*}
 I&=ic^2,&D&=1+\Delta I^2,&Q&=\Delta^2 I^2,&d&=Q-1,\\
 v&=cT,&b&=c+RD,&e&=Q(Dv^2+b^2)-1,&B_a&=2Qvb.
\end{align*}
Notice that $Q=(\Delta I)^2$ is the square needed for rational Pell
saturation; it is not an arbitrary positive auxiliary parameter.

Select one of two pairs $(M,W)$, in both cases putting $L=1+HM$:
\begin{align}
 M_{\rm econ}&=k(I^2+T^2),&
 W_{\rm econ}&=b_0+Ht, \label{eq:economical-new}\\
 C_*&=E+2(X+Y)+sq^2, \nonumber\\
 M_{\rm canc}&=I\bigl(2T+RC_*ic\bigr),&
 W_{\rm canc}&=Lb_0+Ht. \label{eq:cancellation-new}
\end{align}
For a formal variable $x$ let
\begin{align*}
 C(x)&=d(x^2+M^2A)-e,\\
 F_*(x)&=\bigl[C(x)^2+DB_a^2-4d^2x^2M^2A\bigr]^2
                    -4DB_a^2C(x)^2,\\
 F_H(u)&=H^8F_*(u/H),\\
 N&=F_H(W-L\sqrt G)F_H(W+L\sqrt G).
\end{align*}
The last expression is an integer polynomial.  Its radicals describe a
finite conjugate product and are not part of the arithmetic language or
the list of quantified unknowns.  A complete division-free evaluator is
given below.  The product equation and unshifted system are
\begin{equation}
 P=N_kN_t(1+N)-1=0,
 \qquad\text{and}\qquad N_k=1,\quad N_t=1,\quad N=0.
 \label{eq:two-positive-presentations}
\end{equation}

\subsection{A cheaper exact norm evaluator}

The following scalar formulas define $N$ without radicals, division, or
resultants.  They are also the basis of the literal JSON gate lists:
\begin{align*}
 m&=HM,&n&=m^2A,&d_n&=dn,&Z&=H^2Q,\\
 P_H&=ZDv^2,&U_H&=Zb^2,&e_H&=P_H+U_H-H^2,&j&=4P_HU_H,\\
 x_0&=d(W^2+L^2G),&\mathcal T&=2x_0,&
 \kappa&=\bigl[d(W^2-L^2G)\bigr]^2,\\
 \mathcal A&=2(d_n+e_H),&
 \mathcal B&=(d_n-e_H)^2-j,\\
 a_r&=\mathcal T-\mathcal A,&b_r&=\mathcal B-\kappa,&
 c_r&=16d_nj,\\
 v_r&=a_r(a_r\mathcal T+2b_r)-c_r,&
 u_r&=b_r^2-a_r^2\kappa,\\
 N&=u_r(u_r+\mathcal T v_r)+\kappa v_r^2.
\end{align*}
Here $d_n=dn$ means the product of $d$ and $n$; it introduces no new
parameter.  In the emitted circuit $v^2=c^2T^2$, so the separately useful
$T^2$ is computed once.

To prove the evaluator, the elementary identity
\[
 (C^2+j-K)^2-4jC^2=(C^2-K-j)^2-4Kj
\]
rewrites $F_H(u)$, with $x=du^2$, as
\[
 (x^2-\mathcal A x+\mathcal B)^2-c_rx.
\]
For $u=W\pm L\sqrt G$, the two corresponding $x$ values have sum
$\mathcal T$ and product $\kappa$.  Reduction modulo
$x^2-\mathcal T x+\kappa$ first replaces
$x^2-\mathcal A x+\mathcal B$ by $a_rx+b_r$, and then replaces its
square minus $c_rx$ by $u_r+v_rx$.  Its norm is exactly
$u_r^2+\mathcal T u_rv_r+\kappa v_r^2$.  This argument works when the
radicands are squares or zero and when conjugates coincide; no
discriminant is divided out.

Three reassociations each save one multiplication relative to package 30:
compute the common factor $H^2Q$ before forming $P_H,U_H$; compute
$\kappa$ by multiplying before squaring; and use the displayed Horner
form for the final norm.  These are identities over every commutative
ring.  At the same coefficient tuple, the economical polynomial is therefore
unchanged, including all its integer zeros.  A different complete tuple
produces a different fixed polynomial through the same compiler interface.
The sharper construction changes $M,L,W$ as stated;
its new degree is proved below.

\subsection{Weak signs and the inherited interface}

We use three precisely identified consequences of the blank-compiler
theorem.  First, the weak index/transport bootstrap, valid for
$N_k=N_t=\epsilon\in\{1,-1\}$, supplies
\begin{equation}
 \begin{gathered}
 C_w\ge0,\quad q\ge16,\quad X\ge q,\quad Y\ge q^3,\\
 3q<R<q^4-2,\quad E\ge q^4>R+2,\quad
 k=R+\epsilon+hE>E,\quad kY<c<k(Y+1).
 \end{gathered}\label{eq:positive-weak-bootstrap}
\end{equation}
Second, restoring the normalized first, main, strong and auxiliary Pell
equations, together with the main congruence, permits the inherited signed
auxiliary-quotient soundness theorem.  It forces $\epsilon=1$ and extracts
a blank K215 halt.  Third, every blank halt has a canonical positive
normalized witness.  Those canonical witnesses satisfy $C_w>0$ and
$X=2^R$.  These are semantic compiler dependencies, not consequences of
checking four integers modulo a few small primes.

The hypotheses of the weak bootstrap follow from a root of the product
equation before any radical is recovered.  Indeed
\[
 \kappa=x_0^2-4d^2W^2L^2G,
 \qquad N\equiv(u_r+x_0v_r)^2\pmod4.
\]
At a zero of $P$, the three integer factors $N_k,N_t,1+N$ are units.
The last cannot equal $-1$, since this would give $N=-2$, which is not a
square modulo $4$.  Thus $N=0$ and $N_k=N_t=\epsilon$.  For a solution
of the system, $\epsilon=1$ directly.  No consistency assumption is used.

\subsection{All branches and the sharp separation bound}

Factoring the finite norm shows that $N=0$ supplies coherent roots
$g,\tau,f,y$, initially in a finite algebraic extension, with
\begin{equation}
 \begin{gathered}
 g^2=G,\quad \tau^2=A,\quad f^2=D,\\
 dy^2=e-B_af=Q(vf-b)^2-1,\\
 W=Lg+H(M\tau+y).
 \end{gathered}\label{eq:coherent-positive-roots}
\end{equation}
Conversely any such coherent choice makes a factor of $N$ zero.  This
factorization neither divides by a radical nor presupposes the rationality
of any root.  The bounds that follow apply to every complex conjugate of
$y$; such a conjugate can be imaginary.

For the economical weight, put $A_0=a+2$.  From
\eqref{eq:positive-weak-bootstrap},
$A_0<2E$, $c<2kY$, $D<A_0^2I^2<4E^2I^2$, and $c<I^2$.
Every auxiliary branch satisfies
\[
 |y|<2(v\sqrt D+b).
\]
The two elementary estimates
\[
 2v\sqrt D<8kEYIT<\tfrac12M_{\rm econ}kY^2,
 \qquad
 2b<10RE^2I^2<\tfrac12M_{\rm econ}kY^2
\]
use $I^2+T^2\ge2IT$, $kY>8E$, $k>E$, and $Y^2>20R$.
Moreover
\[
 \sqrt A>\ell=EkY>kY^2,
 \qquad \sqrt G-\sqrt A>k(Y^2+Y-1)>kY^2.
\]
For the second inequality use
$\sqrt G>(a+1)c>(E+Y+1)kY$ and
$\sqrt A<\ell+k$.  We obtain the two uniform bounds
\begin{equation}
 |y|<M_{\rm econ}\sqrt A,
 \qquad |M_{\rm econ}\tau+y|<M_{\rm econ}\sqrt G.
 \label{eq:economical-separation}
\end{equation}

For the cancellation weight a tighter, differently shaped estimate is
needed.  Let $S=\Delta I$, so $Q=S^2$.  The weak bootstrap implies
$\Delta>R+2$ and $I^2\ge c^4$, hence
\[
 Q-b=\Delta I^2(\Delta-R)-c-R>0.
\]
Consequently $b\le d$ and $\sqrt D<S<d$.  For
$Z_1=v\sqrt D+b>1$, every auxiliary branch satisfies
\[
 |y|^2\le\frac{QZ_1^2+1}{d}
 <\left[Z_1\left(1+\frac1d\right)\right]^2.
\]
The strict inequality follows on subtracting the left member from the
right: the numerator is $QZ_1^2-d>0$.  Therefore
\begin{equation}
 |y|<Z_1+v+1<cT((a+2)I+1)+c+RD+1.
 \label{eq:sharp-y-new}
\end{equation}
Set $\rho=(1+1/Y)/E$.  Because $c/\ell<\rho$ and
$(a+3)\rho<2$, the $T$ term in \eqref{eq:sharp-y-new}, divided by
$\sqrt A>\ell$, is less than $2TI$.  The remaining large term is
controlled by the exact rational identity
\[
 (a+2)^2\rho
 =E+2Y+X+6+\frac YX+\frac5X+\frac4Y+\frac8E+\frac4{EY}.
\]
Here $Y/X=sq^2/w\le sq^2$ and
$5/X+4/Y+8/E+4/(EY)<1$.  Since $X\ge16$, this gives
$(a+2)^2\rho<C_*-2$.  Also $(c+1)/\ell<2$.  Thus
\[
 \frac{|y|}{\sqrt A}
 <2TI+R(C_*-2)i^2c^3+2
 \le 2TI+RC_*i^2c^3=M_{\rm canc}.
\]
We have proved, without recovering any square root,
\begin{equation}
 |y|<M_{\rm canc}\sqrt A,
 \qquad |M_{\rm canc}\tau+y|<2M_{\rm canc}\sqrt A.
 \label{eq:cancellation-separation}
\end{equation}

Finally, the same weak bounds imply $G>A$ and $b_0>\sqrt A$.  Indeed
$c>kY$ and $\Delta>E^2+2EY$ give
\[
 \Delta c^2>k^2(E^2Y^2+EY)=A-1.
\]
For the second inequality $b_0>akY$ and
\[
 a^2k^2Y^2-A=k^2\{EY(2Y^2-1)+Y^4\}-1>0.
\]

\subsection{Why the smaller carrier is sound}

In the economical construction, the two arguments of $F_*$ associated
with the two main roots are $(W\pm L\sqrt G)/H$.  If $G$ were nonsquare,
the vanishing of one factor of its rational quadratic norm would force
the other to vanish as well.  But the two arguments have distance
$2L\sqrt G/H>2M\sqrt G$, whereas each root of $F_*$ has modulus less
than $M\sqrt G$ by \eqref{eq:economical-separation}.  This contradiction
proves that $G$ is a square.  The selected $g$ is positive: if it were
negative then \eqref{eq:coherent-positive-roots} would give
$W<-L\sqrt G+HM\sqrt G<0$, contrary to $W=b_0+Ht>0$.

For the cancellation construction, the weighted carrier makes the same
conclusion possible with only $L=1+HM$, instead of package 24's
$1+2HM$.  The branch $g=-\sqrt G$ would require
\[
 x=\frac{W+L\sqrt G}{H}
   =t+\frac LH(b_0+\sqrt G)>2M\sqrt A,
\]
because $b_0>\sqrt A$, $\sqrt G>\sqrt A$, and $L/H>M$.
This contradicts \eqref{eq:cancellation-separation}.  Thus the negative
branch is impossible.  If $G$ were nonsquare, rational conjugacy would
force that branch to vanish whenever the positive branch did, again a
contradiction.  Hence $g=\sqrt G$ is a positive integer.

In either construction, work in a finite normal splitting field containing
the roots in \eqref{eq:coherent-positive-roots}.  Now $g$ is rational.
If $\tau$ were irrational, an automorphism would send it to $-\tau$ while
fixing $g$.  Subtracting the two linear relations would require a change
of magnitude $2M\sqrt A$ in $y$.  Every conjugate of $y$ has magnitude
strictly less than $M\sqrt A$, an impossibility.  Thus $\tau$ is rational,
and hence an integer because $\tau^2=A\in\mathbb Z$.
The linear relation makes $y$ rational, and
$f=(e-dy^2)/B_a$ is rational because $B_a>0$.
Since $f^2=D\in\mathbb Z$, $f$ is integral.

\paragraph{Rational Pell saturation.}
If integers $S\ge2,V$ and rational $y$ satisfy
$S^2V^2-(S^2-1)y^2=1$, then $y$ is integral.
For completeness write $y=u/v_0$ in lowest terms, $v_0>0$.
Its equation gives $v_0^2\mid S^2-1$.
The positive nonsquare integer $D_0=(S^2-1)/v_0^2$ has the two
integer-coordinate Pell solutions $(S,v_0)$ and $(S|V|,|u|)$.
Both coordinates in each pair are nonzero.  Write these pairs as
$(\chi_j,\psi_j)$ and $(\chi_l,\psi_l)$ in powers of the fundamental
positive integer-coordinate Pell unit.  The addition formulas give
$\chi_{r+2j}\equiv-\chi_r\pmod{\chi_j}$ and
$\chi_{2j-r}\equiv-\chi_r\pmod{\chi_j}$.
Because $\chi_j\mid\chi_l$, reducing $l$ modulo $2j$ and reflecting
into $[0,j]$ shows that $l$ is an odd multiple of $j$; strict growth
excludes every index smaller than $j$.  The Pell multiplication formulas
then give $\psi_j\mid\psi_l$, so $v_0\mid u$ and $v_0=1$.
This uses the integer-coordinate Pell group, not a larger group of
half-integral units.

Apply the lemma with $S=\Delta I$ and $V=vf-b$.  It proves $y\in\mathbb Z$.
Since $L\equiv1\pmod H$ and in both variants $W\equiv b_0\pmod H$,
the integral linear relation now yields $g\equiv b_0\pmod H$.
Therefore $\gamma=(g-b_0)/H$ is integral.  It is also positive:
$g>(a+1)c$ and $c>k>E>X$ imply
$H\gamma=g-ac-X>c-X>0$.

We have restored exactly
\begin{align*}
 \tau^2-\ell(\ell+k)&=1,\\
 (X+ac+H\gamma)^2-\Delta c^2&=1,\\
 f^2-\Delta i^2c^4&=1,\\
 Q[c(Tf-1)-Rf^2]^2-(Q-1)y^2&=1,
 \qquad N_k=N_t=\epsilon.
\end{align*}
The auxiliary identity uses $f^2=D$ and
$vf-b=c(Tf-1)-Rf^2$.  Absolute values may be taken for $\tau$ and $y$.
If $f<0$, change the signs of $f$ and $T$ together, preserving $Tf$.
The inherited signed-quotient compiler theorem applies, forces
$\epsilon=1$, and recovers a blank K215 halting computation.

Conversely, take a canonical positive normalized compiler witness, with
$g=b_0+H\gamma$ and positive $\tau,f,y$.
For the economical construction choose
\[
 t=\gamma+M(g+\tau)+y>0;
\]
for the cancellation construction choose
\[
 t=L\gamma+M\tau+y>0.
\]
The coherent factor then vanishes and $N_k=N_t=1$.
Every canonical witness can be extended in this way; no new restriction
on its strong Pell index is imposed.
This proves both directions of existential equivalence with the authentic
compiler.  Composing with the K215 syntactic-refutation calibration and
negating existential solvability proves, within ZFC, that absence of
positive integer solutions is equivalent to $\operatorname{Con}(\mathrm{ZFC})$.

\subsection{The shifted systems}

In the system only, replace the transport equation by
\begin{equation}
 (K_0+w)C_w=F+z(q-1).
 \label{eq:shifted-positive-system}
\end{equation}
Here the new positive variable is related to the old one by
$z_{\rm old}=z+1$.  Every shifted solution therefore gives an unshifted
solution, so soundness is immediate.  For completeness use a canonical
compiler witness, for which $C_w>0$, $X=wq=2^R$, $R>3q$, and $F<q$.
Consequently $w>q>F$, and the original transport equation gives
\[
 (z_{\rm old}-1)(q-1)=(K_0+w)C_w-F>0.
\]
Thus $z_{\rm old}\ge2$ and $z=z_{\rm old}-1$ is positive.
This positivity is proved after restoring the compiler, rather than
inserted into the weak bootstrap.  Directly comparing the two sides of
\eqref{eq:shifted-positive-system} saves one addition/subtraction.
The product equations retain the unshifted variable.

\subsection{Exact degrees, including the new cancellation}

All four coefficient ports have degree zero and all twelve supplied
variables have degree one.  The economical norm is the unchanged package
30 polynomial, so its exact degree is $1072$; its product degree is $1081$.
Its leading form, writing $q_0=\beta J$ and $k_0=\eta+\zeta$, is
\[
 (N_{\rm econ})_{1072}
 =2^{32}i^{48}k_0^{96}w^{80}s^{176}q_0^{656}(q_0-F)^{16}.
\]

Here is a full upper and lower degree argument for the new construction.
Replace every supplied variable by $\lambda$ times itself, and use a
subscript $0$ for its highest homogeneous part.  In the rational function
field of the unscaled variables, the relevant square roots have formal
Laurent expansions at $\lambda=\infty$.  Their leading radicands are
nonzero squares, so the ordinary formal binomial expansion constructs all
branches without fractional orders.  The basic degrees are
\[
 \begin{array}{c|rrrrrrrrrr}
  \text{expression}&E,H&Y&c&R&I&d&b&A,G&M&L\\ \hline
  \text{degree}&6&4&5&4&11&46&38&22&27&33.
 \end{array}
\]
Define three nonzero leading forms
\[
 \mathfrak B=R_0E_0^2i^2c_0^4,\qquad
 \mathfrak D=R_0E_0Y_0i^2c_0^4,\qquad
 \mathfrak C=R_0E_0(sq_0^2)i^2c_0^4.
\]
Their degrees are $38,36,35$, respectively.

The first radical satisfies $\sqrt A-\ell=O(\lambda)$.
For every $f$ branch the auxiliary roots satisfy
\[
 y=\varepsilon(b-vf)+O(\lambda^{-8}),\qquad
 \varepsilon\in\{1,-1\},\qquad\deg(vf)=23.
\]
The exact identity
\[
 C_*\ell-\Delta c
 =(2EX+Esq^2-Y^2-4E-4Y-3)c-EC_*\eta
\]
has unique highest homogeneous term $E_0(sq_0^2)c_0$, of degree $14$.
Consequently
\begin{equation}
 M\sqrt A-b=\mathfrak C+\text{terms of order below }35.
 \label{eq:new-cancellation-first}
\end{equation}
Indeed the difference is
$Ri^2c^3(C_*\ell-\Delta c)+O(\lambda^{28})$.

Let $g_+=\sqrt G$ and $\gamma_\pm=(\pm g_+-b_0)/H$.
The small quotient root $\gamma_+$ has order $-1$, hence
$L\gamma_+$ has order $32$.
Also $g_+=ac+O(\lambda^5)$ and
\[
 g_++b_0-2\sqrt A
   =2Y_0c_0+\text{terms of order below }9.
\]
Multiplying by $M$, whose top form is $R_0E_0i^2c_0^3$, gives
\begin{equation}
 M(g_++b_0-2\sqrt A)
   =2\mathfrak D+\text{terms of order below }36.
 \label{eq:new-cancellation-second}
\end{equation}

The exact sixteen-branch factorization is
\begin{equation}
 N=d^8H^{16}\prod(t-y-M\tau-L\gamma),
 \label{eq:new-sixteen-factor}
\end{equation}
where the coherent choices of $g,\tau,f,y$ are all included with
multiplicity.  For $\gamma_+$, four factors have degree $35$ with signed
top forms $\pm\mathfrak C$; the other four have degree $38$ with signed
top forms $\pm2\mathfrak B$.
For $\gamma_-$, the two branches with $\tau=+\sqrt A$ and
$y=+(b-vf)+O(\lambda^{-8})$ have degree $36$ and top form
$2\mathfrak D$, by \eqref{eq:new-cancellation-first} and
\eqref{eq:new-cancellation-second}.
The remaining six have degree $38$: four top forms are $2\mathfrak B$
and two are $4\mathfrak B$.
This exhausts every branch.  It proves the degree upper bound, not just
a degree observed on a numerical specialization.  Multiplying the
nonzero top forms gives
\begin{align*}
 \deg N&=8\cdot46+16\cdot6+4\cdot35+2\cdot36+10\cdot38=1056,\\
 N_{1056}
 &=2^{14}d_0^8H_0^{16}\mathfrak B^{10}\mathfrak C^4\mathfrak D^2\\
 &=2^{46}i^{48}k_0^{96}w^{74}s^{176}q_0^{646}(q_0-F)^{16}.
\end{align*}
This form is nonzero for every fixed $\beta>0$.
Since
\[
 (N_k)_7=-hwsq_0^4,\qquad
 (N_t)_2=w(q_0-F-\alpha)-zq_0,
\]
the product has exact degree $1056+7+2=1065$.
For example its coefficient of
\[
 J^{667}hz\,i^{48}\eta^{96}w^{75}s^{177}
\]
is $2^{46}\beta^{667}>0$.
The three-equation residual degrees are exactly $7,2,1056$; shifting
$z$ by one preserves these leading degrees.

\subsection{Literal cost, checks, and proof status}

The supplied evaluator charges every binary $+,-,\times$, including one
by a fixed numeral.  Intermediate expressions are shared; comparisons
in systems are free under the repository convention.
\begin{center}
\begin{tabular}{lrrrr}
\hline
Construction & Equations & Degree & Multiplications & Total operations\\
\hline
Economical product &1&1081&65&108\\
Economical system &3&1072&63&104\\
Economical shifted system &3&1072&63&103\\
Cancellation product &1&1065&72&118\\
Cancellation system &3&1056&70&114\\
Cancellation shifted system &3&1056&70&113\\
\hline
\end{tabular}
\end{center}
The checks compare each new norm with an independently implemented
eight-dimensional multiquadratic conjugate product on $50$ signed integer
assignments, check the parity and transport identities on those assignments,
and reject dead gates or unused ports.  Full univariate coefficient lists
over two primes confirm the cancellations and the stated leading forms;
this is stronger than merely propagating degree upper bounds through the
gates.  The mathematical arguments above, rather than those finite
examples, establish the unbounded root-recovery theorem and uniform exact
degrees.  No new end-to-end Lean or Rocq check is claimed, and no assertion
that ZFC is consistent is made.

\section{Sixteen unrestricted integers by a cofinal square radix}
\label{sec:signed16}

The new integer construction combines three components that were previously
separate: package 30's compact evaluation of the Jones ratio for its
31-variable quartic; package 29's signed relation combiner; and a new
restriction of the coding radix to a cofinal family of scaled squares.
The last step removes one positivity witness.  A smaller signed separator
also lowers the degree and evaluation cost.  The result is a completely
specified polynomial in \emph{sixteen unrestricted integers}, with
\[
 \deg P_{16}=813056\,5^{32}+10024,
 \qquad \operatorname{ops}(P_{16})=1865.
\]
These are improvements over the delivered 17-integer circuit, not a new
general bound for arbitrary recursively enumerable sets.  In particular,
Sun's 2026 ten-integer theorem is a different, more general representation
theorem, whose compiler is not instantiated here.

\subsection{Fixed index and literal Jones expressions}

Use the literal 31-variable quartic of package 30, supplied in the archive
as inert source data.  Its residuals are quadratic; shifting every positive
coordinate by one and summing their squares gives a normalized polynomial
\(P_{31}(x,\mathbf n)\), independent of the dummy input \(x\).
The source's fifth residual is \(EX+E-aX\), which equals one at the
all-ones tuple, so \(P_{31}(x,\mathbf0)>0\).  Its coefficient of the
fourth power of the retained first Pell root is one; hence its degree is
exactly four.  The quartic is rebound to the \emph{whole} coefficient tuple
proved in the coefficient section; its underlying four-port interface is
unchanged.  This is a semantic use of that interface theorem, not an
inference from matching file hashes.

Put \(\nu=31\) and \(L=5^{32}\).  If \(p_\alpha\) is the coefficient
of the exponent vector \(\alpha\in\mathbb N^{32}\), \(|\alpha|\le4\),
in the shifted quartic, define
\[
 c_\alpha=(4-|\alpha|)!\prod_{r=0}^{31}\alpha_r!\,p_\alpha,
 \qquad e(\alpha)=\sum_{r=0}^{31}\alpha_r5^r.
\]
There are 402 nonzero coefficient records, each a specified polynomial in
the four fixed compiler ports.  The bit lengths of those ports, in order
\((\beta,K_0,M_C,M_F)\), are
\[
 (1064921,1070004,1064921,1064922).
\]
For a coefficient record \(c_\alpha=\sum_t a_t\prod_j C_j^{e_{tj}}\),
the finite exact bound
\[
 |c_\alpha|\le H_\alpha
 :=\sum_t|a_t|2^{\sum_j e_{tj}\operatorname{bits}(C_j)},
 \qquad H_\alpha+1<2^{4269855}
\]
is verified for every record by the supplied index certificate.  Unlisted
records have coefficient zero.  Consequently the following are admissible
fixed integers:
\begin{align*}
 z_*&=2^{4269855}, &
 e_*(T)&=\sum_{|\alpha|\le4}(z_*+c_\alpha)
                          T^{L-e(\alpha)},\\
 y_*&=e_*(2z_*),&
 b_*&=(2z_*)^{L+1}=2^{4269856(L+1)}.
\end{align*}
The corresponding unused Jones index coordinate is
\(u_*=\sum_{r=1}^{31}(2z_*)^{5^r}\).  Base-five uniqueness makes the
exponents in \(e_*\) distinct, and every digit lies strictly between one
and \(2z_*\).  Thus
\[
 2(2z_*)^L\le y_*<(2z_*)^{L+1}=b_*.
\]
The integer \(y_*\) is needed to state the imported index theorem but
\emph{does not occur as a coefficient leaf of the new polynomial}.
It is replaced there by the simple power-of-two coefficient \(b_*\).
None of these fixed integers is an existential parameter.

The eight base variables are
\[
 (e,g_0,h_0,s_0,w_0,\phi,i_0,j_0)\in\mathbb Z^8.
\]
Set \(g=g_0+1\), \(h=h_0+1\), \(s=s_0+1\), \(w=w_0+1\),
\(i=i_0+1\), \(j=j_0+1\), and define
\begin{align*}
 b&=b_*e^2,& \mathcal B&=b^5,&q&=\mathcal B^L,\\
 \ell&=\sum_{r=1}^{31}\mathcal B^{5^r},&
 \lambda&=\sum_{r=0}^{4L-1}\mathcal B^r,&
 N&=q^{16},\\
 \theta&=\mathcal B-2z_*,&
 G&=32b\mathcal B^{5^{31}}.&
\end{align*}
Write \(E_*=e_*(\mathcal B)\), and put
\begin{align*}
 R={}&\bigl[g+E_*q^3+\ell q^5+
 \{2(E_*-z_*\lambda)(1+\mathcal B+g)^4
       +\lambda\mathcal B+\lambda\mathcal Bq^4\}q^7\bigr](N^2-N)\\
 &+\bigl[q^3-b\ell+\ell+\theta\lambda q^3+
              (\mathcal B-2)q^8\bigr](N^2-1).
\end{align*}
Continue with the literal Jones ratio expressions:
\begin{align*}
 U&=N^2w,&Y&=N^2s,&M&=RY,\\
 P&=2M^2U,&K&=R+1+h(P-1),&A&=M(U+1),\\
 B_1&=2R+1,&W_1&=bw,&C&=B_1+W_1+\phi,\\
 J&=4A-5,&Q&=3W_1C-2(W_1^2-1),&D&=(A^2-1)C^2+1,\\
 E&=iJDC^2,&F&=(A^2-1)E^2+1,&G_1&=A+F(CD+1-A),\\
 H&=B_1+jC,&I&=(G_1^2-1)H^2+1.&
\end{align*}
Finally set
\begin{align*}
 a_0&=(P^2-1)K^2+1,&a_1&=DFI,\\
 D_0&=JF,&T_0&=QF+J(H-C),\\
 m&=K^2G-4(C-KY)^2G-8gK^2.&
\end{align*}
Every intermediate symbol is a polynomial abbreviation.  In particular,
all exponents are fixed natural numbers, and no radical or variable
exponent occurs in the delivered arithmetic circuit.

\subsection{The new cofinal-radix lemma}

\paragraph{Lemma (restriction to scaled squares).}
For the fixed admissible Jones index, the source existential predicate
\[
 a_0,a_1\in\operatorname{Sq}(\mathbb Z),\qquad
 D_0\mid T_0,\qquad m>0
 \tag{JR}\label{eq:JR16}
\]
is equivalent to the parent-quartic zero predicate when
\(b=b_*e^2\), \(e\in\mathbb Z\setminus\{0\}\), and the other
seven source coordinates are nonnegative.

\emph{Proof.}
Soundness is immediate from \(b\ge b_*>y_*\): the integer
\(\varepsilon=b-y_*-1\) is nonnegative, and replacing the usual source
definition \(b=y_*+1+\varepsilon\) by this value changes no ratio
expression.  The source soundness theorem applies literally.

For completeness, inspect the constructive proof of
\texttt{Jones1982.USys\_of\_mem} in the pinned
\texttt{Master3.lean}.  For a fixed parent solution \(\mathbf n\),
its choice \(b=2^v\) uses the equality
\(v=x y_*+\sum_r n_r+1\) only to establish the initial bounds
\(v\ge1\), \(x y_*<b\), and \(n_r<b\).
After these bounds, the entire construction uses only those inequalities,
the positive power-of-two condition on \(b\), and the fixed index.
The displayed source was audited for every occurrence of that equality:
there is no later use.  Hence the same ordinary mathematical construction
works for every integer
\(v\ge x y_*+\sum_r n_r+1\).
Write \(b_*=2^\kappa\), where \(\kappa=4269856(L+1)\).
Choose \(t\ge0\) with \(\kappa+2t\) exceeding this bound, take
\(v=\kappa+2t\), and put \(e=2^t\).
Then \(b=2^v=b_*e^2\).  The unchanged
\texttt{Nine.reduced\_of\_mem} construction and
\texttt{Nine.ratio\_complete} give the ratio witnesses.
This is an explicit strengthening of the source construction; no
claim that an arbitrary Diophantine witness may be replaced by a square
is used. \hfill\(\square\)

\subsection{Five base guards and one margin guard}

Only \(\mathbf x=(g_0,s_0,w_0,\phi,j_0)\) needs direct
nonnegativity guards.  The variables \(h_0,i_0\) remain signed.
Introduce six integers \(z_1,\ldots,z_6\).  Put \(C_0=1\) and
\begin{align*}
 d_r&=3x_r+2,& B_r&=1+4d_r C_{r-1}^2z_r^2,&
 C_r&=C_{r-1}B_r\quad(1\le r\le5),\\
 d_6&=3m-1,&B_6&=1+4d_6 C_5^2z_6^2,&C_*&=C_5B_6.
\end{align*}

\paragraph{Lemma (coprime Pell guard).}
If every \(z_r\ne0\), then
\[
 C_*\in\operatorname{Sq}(\mathbb Z)
 \quad\Longrightarrow\quad
 g_0,s_0,w_0,\phi,j_0\ge0,\quad m>0.
\]
Conversely these inequalities permit a choice of nonzero \(z_r\)
for which \(C_*\) is a square.  Moreover, for arbitrary signed
base coordinates and nonzero \(z_r\),
\(C_r\ne0\) and \(|C_*|\ge3^6\).

\emph{Proof.}
Each \(d_r\equiv2\pmod3\) is nonzero, and each
\(B_r\equiv1\pmod4\) is nonzero.  Also
\(\gcd(B_r,C_{r-1})=1\), so the factors \(B_r\) are pairwise
coprime.  If their product is a square, unique factorization makes each
\(|B_r|\) a square.  A negative square is zero or three modulo four;
therefore every \(B_r\) is a positive square.
A negative \(d_r\), with both \(C_{r-1},z_r\ne0\), would give
\(B_r\le-3\).  Thus \(d_r>0\), giving the asserted inequalities.

Conversely, each positive \(d_r\equiv2\pmod3\) is nonsquare.
The Pell equation \(u^2-d_r v^2=1\) has a nontrivial positive solution
whose \(v\) is divisible by any specified positive integer: reduce a
Pell unit matrix modulo that integer and take a positive power equal
to the identity in the finite matrix group.  Take the modulus
\(2|C_{r-1}|\) and set \(z_r=v/(2C_{r-1})\).
This makes \(B_r=u^2\) and \(z_r\ne0\).  Induction proves
completeness.  Finally, without any positivity assumption,
\(|B_r|\ge3\) when \(z_r\ne0\); multiply the six inequalities.
\hfill\(\square\)

\paragraph{Lemma (the two unguarded base coordinates).}
Assume \(e\ne0\), the five base guards are nonnegative,
\(a_0\) is a square, and \(m>0\).  Then \(h\ge1\).
If also \(i\ne0\), replacing \(i\) by \(|i|\) preserves
\eqref{eq:JR16} and makes \(i_0=|i|-1\) nonnegative.

\emph{Proof.}
Here \(b\ge b_*>y_*\), so \(\mathcal B>2z_*\),
\(G>0\), and \(g\ge1\).  From \(m>0\) obtain
\[
 K\ne0,\qquad g<G/8=4b\mathcal B^{5^{31}},\qquad
 4(C-KY)^2<K^2.
\]
Write \(r=5^{31}\).  The first two bounds imply
\(1+\mathcal B+g<\mathcal B^{r+2}\le q\).
Also \(\ell<\mathcal B^{r+1}\), so \(b\ell<q^3\).
In the first bracket defining \(R\), the positive term
\(\lambda\mathcal Bq^4\) strictly absorbs
\(2z_*\lambda(1+\mathcal B+g)^4\); every remaining term is
nonnegative.  The second bracket is positive as well.  Thus
\(R\ge N^2-N\ge2\).
Now \(Y,U\ge1\), \(C=2R+1+bw+\phi>0\), and
\(P=2R^2Y^2U>R+2\).
If \(K<0\), then \(|C-KY|>|K|\), contradicting the displayed
strict inequality.  Hence \(K>0\).
If \(h\le0\), then \(1\le K\le R+1<P\).
The case \(K=1\) would imply \(P-1\mid R\), impossible.
For \(2\le K<P\),
\[
 (PK-1)^2<(P^2-1)K^2+1<(PK)^2,
\]
contradicting squarehood of \(a_0\).  This proves \(h\ge1\).
Finally \(i\) occurs only linearly in \(E\), whose only relevant
consumer is \(E^2\) in \(F\).  Changing the sign of \(i\)
therefore changes none of \(a_0,a_1,D_0,T_0,m\).
\hfill\(\square\)

\subsection{Divisibility and nonvanishing with one extra coordinate}

Use the fixed intersective polynomial
\[
 \Pi(t)=(t^3-t-1)(t^2+23).
\]
It has no rational root.  It has a root modulo every positive integer:
for odd primes other than 23, either the cubic has a simple root, or
it is irreducible and its Frobenius three-cycle makes its discriminant
\(-23\) a nonzero square, giving a simple root of the quadratic.
These roots lift to all prime powers.  At 23, the cubic has the simple
root 3; at 2, \(-23\equiv1\pmod8\) has square roots modulo every
power of two.  The Chinese remainder theorem finishes the proof.

Define the homogeneous degree-five polynomial
\[
 \mathcal R(D,T,v)=
 ((T+Dv)^3-D^2(T+Dv)-D^3)((T+Dv)^2+23D^2).
\]
For integers \(D\ne0,T,Z\),
\[
 (D\mid T\ \text{and}\ Z\ne0)
 \quad\Longleftrightarrow\quad
 \exists v\in\mathbb Z:\ ZD^5\mid\mathcal R(D,T,v).
 \tag{IM}\label{eq:IM16}
\]
Indeed \(\mathcal R=D^5\Pi(T/D+v)\) never vanishes, since
\(T/D+v\) is rational.  If \(D\mid T\), intersectivity supplies
\(v\) so that \(Z\mid\Pi(T/D+v)\).  Conversely, if a prime
satisfies \(\alpha=v_p(D)>\beta=v_p(T)\), then
\(v_p(T+Dv)=\beta\), and the two factors in \(\mathcal R\)
have unique lowest-valuation terms of valuations \(3\beta\)
and \(2\beta\).  Hence
\(v_p(\mathcal R)=5\beta<5\alpha\), contradicting divisibility
by \(D^5\).  The case \(Z=0\) is excluded by nonvanishing.

For the Jones expressions, \(D_0=JF\ne0\) at \emph{every}
integer assignment: \(J=4A-5\ne0\); and if \(F=0\), then
\((1-A^2)E^2=1\) gives \(A=0,E=\pm1\), although
\(E=iJDC^2\) would then be divisible by five.  This contradiction
does not use any guard or desired square condition.

\subsection{The literal sixteen-variable polynomial}

Besides the eight base coordinates and six \(z_r\), introduce
\(v,u\in\mathbb Z\).  These are all sixteen unknowns.  Put
\begin{align*}
 Z&=ie\prod_{r=1}^{6}z_r,& S&=ZD_0^5,&T&=\mathcal R(D_0,T_0,v),\\
 w_1&=C_*,&w_2&=C_*(P^2+K^2+2),&X&=Su-T,\\
 A_0&=S^2C_*,& A_1&=S^2w_1^2a_0,& A_2&=S^2w_2^2a_1,\\
 V&=X^2+A_2-A_0-A_1,&W&=V^2+4X^2A_2-4A_0A_1.&
\end{align*}
The equation is
\[
 \boxed{\quad P_{16}=W^2-16V^2X^2A_2=0.\quad}
\]

\paragraph{Lemma (signed three-root finalizer).}
A zero of this polynomial gives
\[
 S\ne0,\qquad S\mid T,\qquad
 C_*,a_0,a_1\in\operatorname{Sq}(\mathbb Z).
 \tag{F}\label{eq:F16}
\]
Conversely, these conditions on the earlier coordinates allow an
integer choice of \(u\) making the polynomial zero.

\emph{Proof.}
If \(S=0\), then \(A_0=A_1=A_2=0\), \(X=-T\), and
\(P_{16}=T^8\ne0\).  Thus any zero has \(S\ne0\), so
\(i,e,z_1,\ldots,z_6\) are nonzero.  The guard lemma, before
any squarehood conclusion, gives \(|C_*|\ge3^6\).
Consequently
\[
 |w_1|>\sqrt{|C_*|}.
\]
For arbitrary real \(P,K\),
\[
 \sqrt{|(P^2-1)K^2+1|}\le P^2+K^2+1.
\]
To verify this, first bound the absolute radicand by
\((P^2+1)K^2+1\), then subtract this expression from
\((P^2+K^2+1)^2\); the result is
\(P^4+K^4+P^2K^2+2P^2+K^2\ge0\).
It follows that
\[
 |w_2|>\sqrt{|C_*|}+|w_1|\sqrt{|a_0|}.
\]
These estimates permit signed \(C_*,w_1,w_2\); positivity has
not been assumed.

Choose complex square roots \(r_0^2=C_*\), \(r_1^2=a_0\),
\(r_2^2=a_1\).  Direct multiplication gives the polynomial identity
\[
 P_{16}=\prod_{\epsilon_0,\epsilon_1,\epsilon_2\in\{-1,1\}}
 \bigl(X-S(\epsilon_0r_0+\epsilon_1w_1r_1+
                                \epsilon_2w_2r_2)\bigr).
\]
At a zero, one of these signed sums is rational.  Under any
automorphism of the finite Galois extension generated by the three
roots, each root changes only by sign.  If the largest nonzero root
whose sign changes has index 1 or 2, its weighted modulus strictly
exceeds the sum of all earlier weighted moduli by the preceding
bounds: a nonzero root of an integer radicand has modulus at least
one.  A change only at index 0 is impossible as well.  Thus an
automorphism cannot preserve the rational sum while changing a
nonzero root.  Every root is fixed, hence rational, and a rational
square root of an integer is an integer.  This handles negative
candidate radicands, zero roots, and dependent square classes.
The resulting integral signed sum equals \(u-T/S\), so
\(S\mid T\).  This proves necessity.
Conversely, select integral roots and take
\(u=T/S+r_0+w_1r_1+w_2r_2\).  A factor vanishes.
\hfill\(\square\)

\paragraph{Theorem (the concrete integer equation).}
With the fixed coefficients above,
\[
 \mathrm{ZFC}\vdash
 \bigl(\nexists\mathbf t\in\mathbb Z^{16}:
                           P_{16}(\mathbf t)=0\bigr)
 \longleftrightarrow\operatorname{Con}(\mathrm{ZFC}).
\]

\emph{Proof.}
A zero first gives \eqref{eq:F16}.  Apply \eqref{eq:IM16} to
recover \(D_0\mid T_0\) and all nonzero conditions.
The coprime guard gives the five nonnegative base coordinates and
\(m>0\).  The unguarded-coordinate lemma gives \(h_0\ge0\)
and permits replacing \(i\) by \(|i|\).
Because \(e\ne0\), the derived source coordinate
\(\varepsilon=b_*e^2-y_*-1\) is nonnegative.  We have therefore
recovered the literal source Jones predicate, which gives a zero of
the parent quartic and hence, through the whole-tuple compiler and
the ordinary K215 calibration, a finite ZFC refutation.

Conversely, a finite refutation gives a parent-quartic zero.
The cofinal-radix lemma gives Jones witnesses with
\(b=b_*e^2\), \(e\ne0\).  The Pell guard supplies the six
nonzero \(z_r\) and a square \(C_*\).
Intersectivity supplies \(v\) with \(S\mid T\), and the
finalizer supplies \(u\).  All steps are ordinary number theory
and finite Galois theory formalizable in ZFC.  Negating this
two-way existential equivalence proves the assertion.
It proves neither consistency nor inconsistency. \hfill\(\square\)

\subsection{Exact complexity and reproducible checks}

The square-radix substitution doubles the degree contributed by the
radix coordinate, while removing one guard.  With \(L=5^{32}\),
selected exact degrees are
\[
\begin{array}{c|r@{\qquad}c|r}
 R&470L&P&1900L+3\\
 K&1900L+4&m&4442L+12\\
 a_0&7600L+14&a_1&49280L+68\\
 D_0&13750L+20&C_*&4442L+1103\\
 S&68750L+108&T&68750L+105\\
 X&68750L+109&w_2&8242L+1111\\
 A_2&203264L+2506&P_{16}&813056L+10024
\end{array}
\]
The degree of \(A_2\) strictly exceeds those of
\(X^2,A_0,A_1\).  Therefore \(V\) has leading form \(A_2\),
\(W\) has leading form \(A_2^2\), and the leading form of the
final output is \(A_2^4\).  On the diagonal specialization this is
a positive integer times a power of the positive fixed coefficient
\(b_*\).  The supplied leading-term certificate keeps that power
symbolic, establishing the stated exact degree for the authentic
tuple rather than only a generic upper bound.

All fixed powers are expanded into binary multiplications; squaring
and scalar multiplication are charged.  Reuse is free, and fixed
coefficient construction is excluded from variable-evaluation cost.
The emitted, fully live DAG has 705 additions, 26 subtractions,
and 1134 multiplications, totaling 1865.  Its source digit polynomial
uses the compact Horner evaluator and shared radix-five powers of
package 30; no imported Python program is executed.

The independent diagnostics evaluate actual fixed coefficients modulo
three primes and compare 90 signed tuples both with a direct formula
and with the unchanged inherited Jones ratio after substituting
\(\varepsilon=b_*e^2-y_*-1\).  They also check 200 independent
eight-factor norm products on integral roots and another 200 products
in the rank-eight free radical algebra with arbitrary signed radicands,
the basic separator inequalities on
57,222 assignments, and intersectivity modulo every integer through
1200.  These finite tests check definitions and arithmetic.  The
unbounded correctness argument is the proof above, together with
the explicitly identified source compiler and Jones theorems; the
archive does not claim a new end-to-end Lean or Rocq certificate.

\paragraph{Seventeen-integer companion.}
Retaining \(b=y_*+1+\varepsilon\), guarding all six base
coordinates as in package 29, and using the improved separator
\(w_1=C_*\) gives a companion polynomial in seventeen integers.
Its exact degree is \(406528\,5^{32}+27512\), and its 1874 live
operations consist of 709 additions, 26 subtractions, and 1139
multiplications.  It uses the same new coefficient tuple and index.
The corresponding proof is the signed-three-root proof above with
seven guard factors and without the cofinal-square-radix step.
This is an explicit alternative when degree is valued more highly
than the final integer-variable reduction.

\section{Verification, archive contents, and remaining questions}
\label{sec:verification}

\subsection{A reproducible finite audit}

The archive is self-contained for the new finite calculations. Python 3.10
or later and its standard library suffice; no network connection, symbolic
algebra package, proof-assistant installation, or machine simulation is
needed. From the extracted top-level directory run
\begin{lstlisting}
python3 reproduce.py
\end{lstlisting}
The runner authenticates the included source files, rebuilds the four
coefficient integers, runs an independent reconstruction of that tuple,
rebuilds the low-degree and positive circuits, rebuilds both signed
circuits and their fixed-index certificates, and runs the independent
signed audit. It compares the regenerated literal mathematical artifacts
with the delivered versions and writes a receipt under
\texttt{verification/}. The source map includes SHA-256 and pinned Git
blob identities. Imported programs are retained as reference text;
the new builders read their JSON data and do not execute the imported
programs.

The principal checks have the following scopes.
\begin{enumerate}
\item \textbf{Coefficient construction.} The 430 literal actions determine
the 943 ordered tiles. Both implementations check 69230 native-basis
convolution identities, clause-band separation, anchor synchronization,
zero padding fields, carry bounds, population balance, the general
padding controller's orders and valuation, and its unit table. All four
raw byte strings match an independent reconstruction.
\item \textbf{Quadratic and quartic algebra.} An exact sparse polynomial
ring verifies all 22 residual identities after the coordinate shifts,
against both the predecessor's expressions and its literal sparse data.
It also verifies the comparison circuit, the sum of squares, degree four,
the coefficient of $\tau^4$, normalization, and gate liveness.
\item \textbf{Positive norm algebra.} For each retained family, 50 exact
conjugate products in a free radical algebra agree with the scalar
evaluator; 50 parity checks and 50 transport identities also pass.
Full univariate coefficient calculations over each of two prime fields
verify the degree and leading-form predictions. The Laurent argument in
the proof supplies the universal upper bound and nonzero leading form;
the modular calculations are independent implementation diagnostics.
\item \textbf{Signed algebra and fixed index.} Each signed candidate has
an exact symbolic leading-term certificate, live-input/gate checks,
90 modular comparisons at three primes using the actual coefficient
tuple, 200 norm products with integral roots, and 200 further norm
products with arbitrary signed radicands. There are 57222 finite
separator checks and intersectivity checks through modulus 1200.
The separate signed reviewer checks another 300 arbitrary signed-radicand
products and 300 guard stacks, and reads the circuit as inert JSON.
\end{enumerate}

Separate proof reviews examine the infinitary steps that finite arithmetic
cannot settle: the new padding modulus, the sharp positive-root bound,
rational Pell saturation, the cofinal Jones radix, and the order of the
signed nonvanishing and separation arguments. These reviews were performed
by separate agents within this AI-assisted collaboration; they are not
external human endorsement or proof-assistant certification.

\subsection{Which files define the equations}

\begin{center}\small
\begin{tabular}{@{}p{0.31\textwidth}p{0.61\textwidth}@{}}\toprule
Location & Role\\\midrule
\texttt{report.tex}, \texttt{sections/} & Full mathematical exposition and proofs\\
\texttt{report.pdf} & Compiled, visually inspected report\\
\texttt{coefficients/} & Literal machine, four actual integer payloads, finite recipes, binding, and receipts\\
\texttt{positive/} & Six live circuits: two single equations and their ordinary/shifted three-equation systems\\
\texttt{quartic/} & 22-equation quadratic system, quartic sum of squares, and exact sparse data\\
\texttt{integers/generated/} & Sixteen- and seventeen-integer circuits, fixed-index recipes, and degree certificates\\
\texttt{sources/} & Authenticated inherited proof texts, endpoint statements, and source map\\
\texttt{reviews/} & Independent mathematical reviews and separately implemented audits\\
\texttt{RESULTS.json} & Explicit domain, degree, cost, and binding for every delivered form\\
\texttt{MANIFEST.sha256} & File integrity list for the delivered release\\\bottomrule
\end{tabular}
\end{center}

All named intermediate quantities in a circuit are abbreviations for
polynomials, not additional quantified variables. Every fixed coefficient
is either a delivered integer byte string or an exact finite integer
recipe. No coefficient depends on choosing whether ZFC is consistent.
No input port remains an unspecified free parameter in the fixed-instance
bindings.

The two signed polynomials deliberately use package 30's old literal
31-variable quartic, with the complete new coefficient tuple. The new
62-operation system and its 297-term quartic are a separate improvement.
Substituting the latter into the Jones branch would require regenerating
the weighted index coefficients; it is not done implicitly here.

For the PDF, run \texttt{make pdf}, or run \texttt{pdflatex} twice on
\texttt{report.tex}. The supplied Makefile uses no shell escape.
The LaTeX sources use ordinary AMS, Latin Modern, geometry, microtype,
table, listing, and hyperlink packages.

\subsection{Interpretation and next useful reductions}

The new results answer distinct versions of ``simpler.'' The sixteen-integer
construction improves the current explicit signed variable count in this
repository. The twelve-positive-variable cancellation improves degree;
the other twelve-variable presentation improves cost. The quadratic
system gives a short low-degree form, and the coefficient adapter reduces
the common direct coefficient payload. These conclusions do not establish
an optimum in any unrestricted class of Diophantine presentations.

The most consequential unresolved extensions are a fully instantiated
ten-integer construction based on Sun's different representation, further
removal of signed positivity witnesses, and a kernel-checked proof of the
new compiler adapter and finalizers. An improvement to the nine-positive-
variable count or to the established 67-operation equation would require
another idea; this report does not claim those improvements.

The result proves an equivalence. Assuming ZFC is consistent, G\"odel's
second incompleteness theorem prevents ZFC from proving that the equation
has no solutions. ZFC can nevertheless prove that each solution encodes
a finite refutation, and that every finite refutation supplies a solution.


\begin{thebibliography}{99}
\bibitem{repo} Vladimir Reshetnikov, \emph{ProveIt}, frozen commit
\pin. Repository: \repo. The bundled source map gives the precise path
of each imported text, literal data file, and construction.
\bibitem{p24} ProveIt package \texttt{24-integer24},
\emph{Smaller explicit Diophantine forms of ZFC consistency}.
Especially \texttt{sections/coefficients.tex},
\texttt{sections/positive.tex}, and the coefficient generator.
\bibitem{p25} ProveIt package \texttt{25-hybrid12},
\emph{Smaller Diophantine forms of ZFC consistency}.
Especially \texttt{coefficients/PROOF.md},
\texttt{proofs/HYBRID12\_PROOF.md}, and the inherited compiler notes.
\bibitem{p27} ProveIt package \texttt{27-wave4},
\emph{Smaller Diophantine presentations of ZFC consistency}.
Especially the Wang943 compiler/interface proofs and the retained
transport-shift proof.
\bibitem{p29} ProveIt package \texttt{29-simpler-forms},
\emph{Smaller explicit Diophantine forms of ZFC consistency}.
Especially \texttt{signed17/SIGNED17\_PROOF.md}.
\bibitem{p30} ProveIt package \texttt{30-quartic31},
\emph{Smaller explicit Diophantine forms for Con(ZFC)}.
Especially the literal quartic, compact Jones evaluator, and
\texttt{positive/POSITIVE12\_PROOF.md}.
\bibitem{jones} James P. Jones, \emph{Universal Diophantine equation},
Journal of Symbolic Logic 47 (1982), 549--571.
\url{https://doi.org/10.2307/2273588}.
\bibitem{sun10} Zhi-Wei Sun, \emph{Ten unknowns for Hilbert's tenth problem
over the integers}, arXiv:2609.01594v2, 2 September 2026, Theorem 1.1.
\url{https://arxiv.org/abs/2609.01594v2}.
\end{thebibliography}
\end{document}
